\chapter{Théorème Central Limite}
\sloppy


\section{Le théorème de Paul Lévy}
\begin{theorem}[Paul Lévy]
Soit $(X_n)_{n \in \mathbb{N}}$ une suite de v.a. réelles et soit $X$ une v.a. réelle. Soit $\phi_{X_n}$, $\phi_X$ leurs fonctions caractéristiques. Nous avons équivalence entre:
\begin{enumerate}
    \item[i)] $X_n \xrightarrow{\mathcal{L}} X$
    \item[ii)] $\forall u \in \mathbb{R}, \lim_{n \to \infty} \phi_{X_n}(u) = \phi_X(u)$
\end{enumerate}
\end{theorem}
\begin{proof}
Sens facile : i) $\implies$ ii) \\
On a $\phi_{X_n}(u) = \mathbb{E}[e^{iuX_n}] = \mathbb{E}[\cos(uX_n)] + i\mathbb{E}[\sin(uX_n)]$.
Les fonctions $x \mapsto \cos(ux)$ et $x \mapsto \sin(ux)$ (u étant fixé) sont continues et bornées. Par définition de la convergence en loi,
\begin{align*}
    \mathbb{E}[\cos(uX_n)] &\xrightarrow[n \to \infty]{} \mathbb{E}[\cos(uX)] \\
    \mathbb{E}[\sin(uX_n)] &\xrightarrow[n \to \infty]{} \mathbb{E}[\sin(uX)]
\end{align*}
et donc $\phi_{X_n}(u) \xrightarrow[n \to \infty]{} \phi_X(u)$.
\vspace{1cm}
Sens réciproque : ii) $\implies$ i) \\
Nous utilisons la même stratégie que pour prouver l'injectivité de la transformée de Fourier des mesures de probabilités. Nous supposons que $X, X_n$ pour $n \in \mathbb{N}$ sont toutes définies sur un même espace de probas $(\Omega, \mathcal{F}, P)$, sur lequel est aussi définie une v.a. $Y$ indépendante de $X, X_n$ pour $n \in \mathbb{N}$, et de la loi $\mathcal{N}(0,1)$.
Soit $f: \mathbb{R} \to \mathbb{R}$ une fonction continue à support compact. Nous supposons ii) vrai, et nous voulons montrer que $X_n \xrightarrow{\mathcal{L}} X$. Soit $\sigma > 0$ et écrivons :
\begin{align*}
    \mathbb{E}[f(X_n)] - \mathbb{E}[f(X)] &= \mathbb{E}[f(X_n) - f(X_n + \sigma Y)] + \mathbb{E}[f(X_n + \sigma Y)] - \mathbb{E}[f(X + \sigma Y)] \\
    &+ \mathbb{E}[f(X + \sigma Y) - f(X)]
\end{align*}
Comme $f$ est continue à support compact, elle est uniformément continue.
\[
\forall \varepsilon > 0, \exists \delta > 0, \forall x,y \in \mathbb{R}, |x-y| < \delta \implies |f(x) - f(y)| < \varepsilon
\]
Soit $\varepsilon > 0$ fixé. Soit $\delta$ associé à $\varepsilon$ comme ci-dessus. Écrivons :
\begin{align*}
    |\mathbb{E}[f(X_n) - f(X_n + \sigma Y)]| &\leq \mathbb{E}[|f(X_n) - f(X_n + \sigma Y)|] \\
    &= \int_{|\sigma Y| \leq \delta} |f(X_n) - f(X_n + \sigma Y)| dP + \int_{|\sigma Y| > \delta} |f(X_n) - f(X_n + \sigma Y)| dP
\end{align*}
\begin{align*}
    \int_{|\sigma Y| \leq \delta} |f(X_n) - f(X_n + \sigma Y)| dP &< \varepsilon \\
    \int_{|\sigma Y| > \delta} |f(X_n) - f(X_n + \sigma Y)| dP &\leq \int_{|\sigma Y| > \delta} 2\|f\|_{\infty} dP = 2\|f\|_{\infty} \mathbb{P}(|\sigma Y| > \delta)
\end{align*}
On a $\lim_{\sigma \to 0} \mathbb{P}(|\sigma Y| > \delta) = \lim_{\sigma \to 0} \mathbb{P}(|Y| > \delta/\sigma) = 0$. Donc $\exists \sigma_0$ tel que pour $\sigma < \sigma_0$, $2\|f\|_{\infty} \mathbb{P}(|\sigma Y| > \delta) < \varepsilon$.
On a donc $|\mathbb{E}[f(X_n) - f(X_n + \sigma Y)]| < 2\varepsilon$.
Par le même argument, on a $|\mathbb{E}[f(X + \sigma Y) - f(X)]| < 2\varepsilon$.
Reste le terme du milieu.
\[
\forall n \in \mathbb{N}, \quad |\mathbb{E}[f(X_n)] - \mathbb{E}[f(X)]| \leq 2\varepsilon + |\mathbb{E}[f(X_n + \sigma Y) - f(X + \sigma Y)]| + 2\varepsilon
\]
Nous utilisons la formule du cours précédent (via la convolution avec une Gaussienne) :
\[
\mathbb{E}[f(X_n + \sigma Y)] = \iint f(z) e^{iuz} \frac{1}{\sqrt{2\pi}\sigma} g_{1/\sigma}(u) \phi_{X_n}(-u) du dz
\]
\[
\mathbb{E}[f(X + \sigma Y)] = \iint f(z) e^{iuz} \frac{1}{\sqrt{2\pi}\sigma} g_{1/\sigma}(u) \phi_{X}(-u) du dz
\]
où $g_{1/\sigma}$ est la densité de $\mathcal{N}(0, 1/\sigma^2)$.
Posons $F_n(u,z) = f(z) e^{iuz} \frac{1}{\sqrt{2\pi}\sigma} g_{1/\sigma}(u) \phi_{X_n}(-u)$. On a :
\[
|F_n(u,z)| \leq |f(z)| \frac{1}{\sqrt{2\pi}\sigma} g_{1/\sigma}(u)
\]
Cette majoration ne dépend pas de $n$ et est dans $L^1(d\lambda(u), d\lambda(z)) = L^1(\mathbb{R}^2, \lambda^{\otimes 2})$.
Par le théorème de convergence dominée (TCD), comme $\phi_{X_n}(u) \to \phi_X(u)$, on a :
\[
\exists N, \forall n \geq N, \quad |\mathbb{E}[f(X_n + \sigma Y)] - \mathbb{E}[f(X + \sigma Y)]| < \varepsilon
\]
Alors, pour $n \geq N$,
\[
|\mathbb{E}[f(X_n)] - \mathbb{E}[f(X)]| < 2\varepsilon + \varepsilon + 2\varepsilon = 5\varepsilon
\]
Ceci étant vrai pour tout $\varepsilon > 0$, on a $\mathbb{E}[f(X_n)] \to \mathbb{E}[f(X)]$ pour toute fonction $f$ continue à support compact, ce qui est la définition de la convergence en loi.
\end{proof}
\section{Moments}
\begin{proposition}
Soit $X$ une v.a. réelle. Soit $k \geq 1$. Supposons que $\mathbb{E}[|X|^k] < \infty$.
Alors $\phi_X$ est $k$ fois dérivable, et de plus
\[ \phi_X^{(k)}(0) = i^k \mathbb{E}[X^k] \]
\end{proposition}
\begin{proof}
Faisons la preuve pour $k=1$. On a $\mathbb{E}[|X|] < \infty$.
La fonction caractéristique est $\phi_X(u) = \int_{\mathbb{R}} e^{iux} dP_X(x)$.
Il s'agit d'une intégrale à paramètre, et on veut dériver par rapport à $u$.
Posons $F(x,u) = e^{iux}$. On a $\frac{\partial F}{\partial u}(x,u) = ix e^{iux}$.
On a $\left|\frac{\partial F}{\partial u}(x,u)\right| = |ix e^{iux}| = |x|$.
La fonction majorante $|x|$ est $P_X$-intégrable car $\mathbb{E}[|X|] < \infty$, et elle ne dépend pas de $u$.
Par le théorème de dérivation sous le signe intégral, $\phi_X$ est dérivable et
\[ \phi'_X(u) = \int_{\mathbb{R}} \frac{\partial F}{\partial u}(x,u) dP_X(x) = \int_{\mathbb{R}} ix e^{iux} dP_X(x) \]
D'où $\phi'_X(0) = \int_{\mathbb{R}} ix dP_X(x) = i \mathbb{E}[X]$.
Le résultat pour $k > 1$ s'obtient par récurrence.
\end{proof}
Ce résultat permet de trouver facilement la suite des moments.
\subsection*{Exemples}
\begin{example}[Loi uniforme]
Soit $X \sim \mathcal{U}([-1,1])$. La densité est $f(x) = \frac{1}{2} \mathbf{1}_{[-1,1]}(x)$.
\begin{align*}
    \phi_X(u) = \int_{-1}^1 e^{iux} \frac{1}{2} dx = \frac{1}{2} \left[ \frac{e^{iux}}{iu} \right]_{-1}^1 = \frac{e^{iu} - e^{-iu}}{2iu} = \frac{\sin(u)}{u}
\end{align*}
On connaît le développement en série de $\sin(u) = \sum_{n=0}^\infty (-1)^n \frac{u^{2n+1}}{(2n+1)!}$.
D'où $\phi_X(u) = \frac{\sin(u)}{u} = \sum_{n=0}^\infty (-1)^n \frac{u^{2n}}{(2n+1)!}$.
La dérivée d'ordre $2n$ en $0$ est $(2n)!$ fois le coefficient de $u^{2n}$:
\[ \phi_X^{(2n)}(0) = (2n)! \frac{(-1)^n}{(2n+1)!} \]
D'où, $\mathbb{E}[X^{2n}] = \frac{\phi_X^{(2n)}(0)}{i^{2n}} = \frac{1}{(-1)^n} (2n)! \frac{(-1)^n}{(2n+1)!} = \frac{1}{2n+1}$.
Plus facile: $\mathbb{E}[X^{2n}] = \int_{-1}^1 x^{2n} \frac{1}{2} dx = \frac{1}{2} \left[ \frac{x^{2n+1}}{2n+1} \right]_{-1}^1 = \frac{1}{2n+1}$.
\end{example}
\begin{example}[Loi Gaussienne]
Soit $X \sim \mathcal{N}(0,1)$. On sait que $\phi_X(u) = e^{-u^2/2}$.
Le développement en série est :
\[ \phi_X(u) = e^{-u^2/2} = \sum_{n=0}^\infty \frac{1}{n!} \left( -\frac{u^2}{2} \right)^n = \sum_{n=0}^\infty \frac{(-1)^n}{n! 2^n} u^{2n} \]
On en déduit que $\phi_X^{(2n)}(0) = (2n)! \frac{(-1)^n}{n! 2^n}$.
D'où $\mathbb{E}[X^{2n}] = \frac{\phi_X^{(2n)}(0)}{i^{2n}} = \frac{1}{(-1)^n} \frac{(2n)! (-1)^n}{n! 2^n} = \frac{(2n)!}{n! 2^n}$.
\end{example}
\section{Convolution}
\begin{lemma}
Soient $X, Y$ deux v.a. indépendantes. Alors $\forall u \in \mathbb{R}$,
\[ \phi_{X+Y}(u) = \phi_X(u) \phi_Y(u) \]
\end{lemma}
\begin{proof}
$\phi_{X+Y}(u) = \mathbb{E}[e^{iu(X+Y)}] = \mathbb{E}[e^{iuX} e^{iuY}]$. Par indépendance de $X$ et $Y$, les v.a. $e^{iuX}$ et $e^{iuY}$ sont aussi indépendantes. Donc,
\[ \mathbb{E}[e^{iuX} e^{iuY}] = \mathbb{E}[e^{iuX}] \mathbb{E}[e^{iuY}] = \phi_X(u) \phi_Y(u) \]
\end{proof}
\begin{corollary}
Si $X_1, \dots, X_n$ sont $n$ v.a. indépendantes, alors $\phi_{X_1+\dots+X_n}(u) = \prod_{i=1}^n \phi_{X_i}(u)$.
Si de plus elles sont identiquement distribuées, $\phi_{X_1+\dots+X_n}(u) = (\phi_{X_1}(u))^n$.
\end{corollary}
Supposons que $X$ a une loi de densité $f$ et $Y$ une loi de densité $g$, et $X,Y$ indépendantes.
\[ \phi_{X+Y}(u) = \iint_{\mathbb{R}^2} e^{iu(x+y)} f(x)g(y) dx dy \]
On pose $z=x+y$. On garde $x$ et on remplace $y=z-x$.
\begin{align*}
    \phi_{X+Y}(u) &= \iint e^{iuz} f(x)g(z-x) dx dz \quad \text{(Fubini)} \\
    &= \int e^{iuz} \left( \int f(x)g(z-x) dx \right) dz = \int e^{iuz} h(z) dz
\end{align*}
où $h(z) = \int f(x)g(z-x) dx = (f*g)(z)$ est le produit de convolution de $f$ et $g$.
Par identification de ces deux formules, on voit que $X+Y$ a une loi de densité $f*g$.
Dès que $f$ et $g$ sont dans $L^1$, le produit de convolution $f*g$ est bien défini et aussi dans $L^1$.
\section{Théorème Central Limite}
\begin{theorem}[Théorème Central Limite]
Soit $(X_n)_{n \geq 1}$ une suite de v.a. i.i.d. (indépendantes et identiquement distribuées) qui admet un moment d'ordre 2. Posons $m = \mathbb{E}[X_i]$ et $\sigma^2 = \text{Var}[X_i]$. Alors,
\[ \frac{X_1 + \dots + X_n - nm}{\sqrt{n}} \xrightarrow[n \to \infty]{\mathcal{L}} \mathcal{N}(0, \sigma^2) \]
\end{theorem}
\textbf{Attention :} Il faut absolument centrer les v.a. en soustrayant $\mathbb{E}[X_1+\dots+X_n] = nm$. \\
Par contre, on peut normaliser pour obtenir une variance de 1 :
\[ \frac{X_1 + \dots + X_n - nm}{\sqrt{n}\sigma} \xrightarrow[n \to \infty]{\mathcal{L}} \mathcal{N}(0, 1) \]
\subsection*{Histoire}
Il s'agit d'un résultat universel. Dès que les v.a. $(X_n)$ ont un moment d'ordre 2, les fluctuations de $X_1, \dots, X_n$ autour de sa moyenne sont décrites par la loi gaussienne. Le TLC (ou TCL) est un raffinement de la loi des grands nombres, sous l'hypothèse d'existence d'un moment d'ordre 2.
\begin{itemize}
    \item \textbf{Abraham de Moivre (1733)} : \textit{The doctrine of chances}.
    \item \textbf{Pierre Simon de Laplace (1812)} : \textit{Traité de Mécanique céleste} et \textit{Théorie analytique des probabilités} (CLT pour la loi binomiale, théorème de Moivre-Laplace).
    \item \textbf{Lindeberg}, \textbf{Turing (1934)}.
\end{itemize}
\begin{proof}[Preuve du TCL]
On pose $Y_n = \frac{X_1 + \dots + X_n - nm}{\sqrt{n}}$ et $\tilde{X}_k = X_k - m$ pour $1 \leq k \leq n$.
On a $\mathbb{E}[\tilde{X}_k] = 0$ et $\mathbb{E}[\tilde{X}_k^2] = \sigma^2$.
On peut réécrire $Y_n = \frac{1}{\sqrt{n}} \sum_{k=1}^n \tilde{X}_k$.
Étudions la convergence en loi de $Y_n$ à l'aide des fonctions caractéristiques. Les $\tilde{X}_k$ sont i.i.d.
\begin{align*}
    \phi_{Y_n}(u) = \mathbb{E}\left[ e^{iu Y_n} \right] = \mathbb{E}\left[ e^{i \frac{u}{\sqrt{n}} \sum_{k=1}^n \tilde{X}_k} \right] &= \mathbb{E}\left[ \prod_{k=1}^n e^{i \frac{u}{\sqrt{n}} \tilde{X}_k} \right] \\
    &= \prod_{k=1}^n \mathbb{E}\left[ e^{i \frac{u}{\sqrt{n}} \tilde{X}_k} \right] \quad \text{(par indépendance)} \\
    &= \left( \phi_{\tilde{X}_1}\left(\frac{u}{\sqrt{n}}\right) \right)^n \quad \text{(car i.i.d.)}
\end{align*}
Comme $\mathbb{E}[\tilde{X}_1^2] < \infty$, $\phi_{\tilde{X}_1}$ est deux fois dérivable. Elle admet donc un développement limité à l'ordre 2 en 0 :
\[ \phi_{\tilde{X}_1}(v) = \phi_{\tilde{X}_1}(0) + v \phi'_{\tilde{X}_1}(0) + \frac{v^2}{2} \phi''_{\tilde{X}_1}(0) + o(v^2) \]
Calculons les coefficients:
\begin{itemize}
    \item $\phi_{\tilde{X}_1}(0) = \mathbb{E}[e^0] = 1$
    \item $\phi'_{\tilde{X}_1}(0) = i \mathbb{E}[\tilde{X}_1] = i \cdot 0 = 0$
    \item $\phi''_{\tilde{X}_1}(0) = i^2 \mathbb{E}[\tilde{X}_1^2] = - \sigma^2$
\end{itemize}
Donc, $\phi_{\tilde{X}_1}(v) = 1 - \frac{\sigma^2 v^2}{2} + o(v^2)$.
Appliquons ce développement avec $v = u/\sqrt{n}$:
\[ \phi_{\tilde{X}_1}\left(\frac{u}{\sqrt{n}}\right) = 1 - \frac{\sigma^2}{2} \left(\frac{u}{\sqrt{n}}\right)^2 + o\left(\left(\frac{u}{\sqrt{n}}\right)^2\right) = 1 - \frac{\sigma^2 u^2}{2n} + o\left(\frac{1}{n}\right) \]
Maintenant, on élève à la puissance $n$:
\[ \phi_{Y_n}(u) = \left( 1 - \frac{\sigma^2 u^2}{2n} + o\left(\frac{1}{n}\right) \right)^n \]
Pour trouver la limite, on utilise le logarithme (complexe). On sait que pour $z \to 0$, $\ln(1+z) = z + O(z^2)$.
\begin{align*}
    \ln(\phi_{Y_n}(u)) &= n \ln\left( 1 - \frac{\sigma^2 u^2}{2n} + o\left(\frac{1}{n}\right) \right) \\
    &= n \left( -\frac{\sigma^2 u^2}{2n} + o\left(\frac{1}{n}\right) \right) \\
    &= -\frac{\sigma^2 u^2}{2} + n \cdot o\left(\frac{1}{n}\right) = -\frac{\sigma^2 u^2}{2} + o(1)
\end{align*}
Ainsi, $\lim_{n \to \infty} \ln(\phi_{Y_n}(u)) = -\frac{\sigma^2 u^2}{2}$.
Par continuité de l'exponentielle,
\[ \lim_{n \to \infty} \phi_{Y_n}(u) = e^{-\frac{\sigma^2 u^2}{2}} \]
On reconnaît la fonction caractéristique de la loi $\mathcal{N}(0, \sigma^2)$.
On conclut par le théorème de Paul Lévy que $Y_n \xrightarrow{\mathcal{L}} \mathcal{N}(0, \sigma^2)$.
\end{proof}
